\chapter{复数}\label{ch:b1:complex}

复数在高中分册中曾作为解二次方程的计算工具出现过。
本章则把复数作为核心对象来研究：
指数形式及其推论（棣莫弗公式、$n$ 次方根、单位根），
复数域 $\C$ 与平面几何之间的系统对应，
以及 $\eu^{\iu\theta}$ 作为
证明三角恒等式的工具。

\section{复数域 $\C$、模与共轭复数}

\begin{definition}[复数域]\label{def:b1:complex:field}
$\C = \{a + \iu b : a, b \in \R\}$，其上通常的加法与
由 $\iu^2 = -1$ 确定的乘法。每个非零的 $z = a + \iu
b$ 有逆元 $z^{-1} = \frac{a - \iu b}{a^2 + b^2}$ --- 用\cref{ch:b1:structures}中的语言说，
$\C$ 是一个域。记 $a =
\Re(z)$，$b = \Im(z)$，$\conj{z} = a - \iu b$ 称为
\emph{共轭复数}\index{共轭复数}，$\abs{z} = \sqrt{a^2 + b^2}$ 称为
\emph{模}\index{模}。
\end{definition}

\begin{proposition}[共轭与模的运算规则]\label{prop:b1:complex:rules}
设 $z, w \in \C$，则：
\begin{enumerate}
  \item $\conj{z + w} = \conj z + \conj w$, $\conj{zw} = \conj z\,
        \conj w$, $\conj{\conj z} = z$;
  \item $z \conj z = \abs z^2$; $\;\Re(z) = \frac{z + \conj z}{2}$,
        $\Im(z) = \frac{z - \conj z}{2\iu}$;
  \item $\abs{zw} = \abs z\, \abs w$，以及 $\abs{z^{-1}} =
        \abs{z}^{-1}$，其中 $z \neq 0$；
  \item （三角不等式\index{三角不等式}）$\abs{z + w}
        \leq \abs z + \abs w$，等号成立当且仅当 $z$ 与
        $w$ 位于从原点 $0$ 出发的同一条射线上（$w = \lambda z$ 或
        $z = \lambda w$，其中 $\lambda \geq 0$）；
  \item （反向三角不等式）$\bigl|\abs z - \abs w\bigr|
        \leq \abs{z - w}$。
\end{enumerate}
\end{proposition}

\begin{proof}
(1)与(2)是对实部与虚部的直接计算。(3)：
$\abs{zw}^2 = zw\,\conj{zw} = z\conj z\, w \conj w = \abs z^2 \abs
w^2$，开平方即得结论；对 $z \cdot z^{-1} = 1$ 取模，
即得逆元的情形。

(4) 两边均非负，故比较平方：
\[
\abs{z+w}^2 = (z+w)(\conj z + \conj w)
= \abs z^2 + \abs w^2 + 2\,\Re(z \conj w),
\]
且 $\Re(z\conj w) \leq \abs{z \conj w} = \abs z \abs w$，由此
$\abs{z+w}^2 \leq (\abs z + \abs w)^2$。等号成立时必有 $\Re(z \conj
w) = \abs{z \conj w}$，即 $z \conj w \in \R_+$；若 $w \neq 0$，
则 $z = \frac{z\conj w}{\abs w^2}\, w = \lambda w$，其中 $\lambda
\geq 0$（$w = 0$ 的情形显然）。

(5) 由 $\abs z = \abs{(z - w) + w} \leq \abs{z-w} + \abs w$ 得 $\abs z
- \abs w \leq \abs{z - w}$；交换 $z$ 与 $w$ 即得另一侧不等式。
\end{proof}

\begin{example}[模与辐角的完整演算]\label{ex:b1:complex:workout}
将 $w = \dfrac{3 + 4\iu}{1 - 2\iu}$ 化为代数形式，
并用两种方法求其模。分子分母同乘以
分母的共轭复数：
\[
w = \frac{(3 + 4\iu)(1 + 2\iu)}{(1 - 2\iu)(1 + 2\iu)}
= \frac{3 + 6\iu + 4\iu - 8}{1 + 4}
= \frac{-5 + 10\iu}{5} = -1 + 2\iu .
\]
直接计算得 $\abs w = \sqrt{1 + 4} = \sqrt5$。利用商的性质
（\cref{prop:b1:complex:rules}~(3)）：$\abs w = \frac{\abs{3 +
4\iu}}{\abs{1 - 2\iu}} = \frac5{\sqrt5} = \sqrt5$ --- 结果相同，
无需代数形式。这一规律具有一般性：模与辐角便于处理积与商，
而实部与虚部便于处理和。
因此应选择与所作运算相匹配的表示，
并只在
必要时才作转换。
\end{example}

\begin{example}[涉及共轭复数的方程]\label{ex:b1:complex:conjeq}
在 $\C$ 中求解 $\;z + 2\conj z = 6 + 2\iu$。同时含 $z$ 与
$\conj z$ 的方程\emph{不是}关于 $z$ 的多项式；可靠的做法是
把它拆成实部与虚部两部分。记 $z = x + \iu y$：
\[
z + 2\conj z = 3x - \iu y ,
\]
于是方程化为 $3x = 6$ 与 $-y = 2$，唯一解为
$z = 2 - 2\iu$。（验证：$(2 - 2\iu) + 2(2 + 2\iu) = 6 +
2\iu$。）另一种做法是对整个方程取共轭，得到
$\conj z + 2z = 6 - 2\iu$，然后解以
$z, \conj z$ 为未知数的线性方程组 --- 结果相同，这在系数为复数时是一个有用的技巧。
关于 $z$ 与 $\conj z$ 的方程
实际上是两个实方程构成的方程组；断言“次数为 $1$、
恰有一解”在此是安全的，但 $z\conj z = -1$（无解）表明，
一旦出现乘积，多项式的直觉便会失效。
\end{example}

\section{指数形式}

\begin{definition}[纯虚指数的复指数]\label{def:b1:complex:expi}
对 $\theta \in \R$，规定
\[
\eu^{\iu\theta} = \cos\theta + \iu \sin\theta .
\]
每个 $z \neq 0$ 都可写成 $z = r\,\eu^{\iu\theta}$，其中 $r =
\abs z > 0$；$\theta$ 称为 $z$ 的\emph{辐角}\index{辐角}，
在相差 $2\pi$ 的整数倍的意义下确定。位于
$\intoc{-\pi}{\pi}$ 中的那一个是\emph{主辐角}，记作
$\arg z$。
\end{definition}

\begin{example}[若干初值与一个著名恒等式]\label{ex:b1:complex:eulerid}
在几个特殊角上直接由定义读出：
\[
\eu^{\iu\pi/2} = \iu, \qquad
\eu^{\iu\pi} = -1, \qquad
\eu^{2\iu\pi} = 1, \qquad
\eu^{\iu\pi/4} = \frac{\sqrt2}2\,(1 + \iu) .
\]
第二个式子移项得 $\eu^{\iu\pi} + 1 = 0$，即为欧拉
著名的恒等式，它联系了 $\eu$、$\iu$、$\pi$、$1$ 与 $0$。
在本阶段它只是定义的直接展开，而非定理，其真正含义 ---
为何将 \cref{ch:b1:functions} 中的解析指数函数
推广到虚变量后仍沿用同一名称 --- 由\cref{ch:b1:series}
中的幂级数予以确定。
与此同时，上面的式子
值得熟记，将它当作一张换算表使用，
每当从图上读取辐角时都会默默用到它。
\end{example}

\begin{theorem}[函数方程]\label{thm:b1:complex:funceq}
对所有 $\theta, \varphi \in \R$：
\[
\eu^{\iu\theta}\, \eu^{\iu\varphi} = \eu^{\iu(\theta + \varphi)},
\qquad
\abs{\eu^{\iu\theta}} = 1,
\qquad
\conj{\eu^{\iu\theta}} = \eu^{-\iu\theta} = (\eu^{\iu\theta})^{-1}.
\]
由此，$\abs{zw} = \abs z \abs w$ 对辐角同样成立，即
$\arg(zw) \equiv \arg z + \arg w \pmod{2\pi}$.
\end{theorem}

\begin{proof}
展开乘积，并利用加法公式：
\[
(\cos\theta + \iu\sin\theta)(\cos\varphi + \iu\sin\varphi)
= (\cos\theta\cos\varphi - \sin\theta\sin\varphi)
+ \iu\,(\sin\theta\cos\varphi + \cos\theta\sin\varphi),
\]
即 $\cos(\theta+\varphi) + \iu\sin(\theta+\varphi)$。模
为 $\sqrt{\cos^2\theta + \sin^2\theta} = 1$，共轭公式
由余弦与正弦的奇偶性得到，它是 $\eu^{\iu\theta}$ 的逆，
因为 $\eu^{\iu\theta}\eu^{-\iu\theta} = \eu^0 = 1$。
\end{proof}

\begin{corollary}[棣莫弗公式]\label{cor:b1:complex:demoivre}
对 $\theta \in \R$ 与 $n \in \Z$：
$\;(\cos\theta + \iu\sin\theta)^n = \cos(n\theta) +
\iu\sin(n\theta)$。\index{棣莫弗公式}
\end{corollary}

\begin{proof}
当 $n \geq 0$ 时用归纳法。$n = 0$ 时为 $1 = 1$。假设
公式对 $n$ 成立，由函数方程
（\cref{thm:b1:complex:funceq}）可得
\[
\bigl(\eu^{\iu\theta}\bigr)^{n+1}
= \bigl(\eu^{\iu\theta}\bigr)^{n}\,\eu^{\iu\theta}
= \eu^{\iu n\theta}\,\eu^{\iu\theta}
= \eu^{\iu(n+1)\theta} ,
\]
即为 $n + 1$ 时的公式。对 $n < 0$，记 $n = -m$
其中 $m > 0$：由于 $\eu^{\iu m\theta}$ 的逆为
$\eu^{-\iu m\theta}$（再由 \cref{thm:b1:complex:funceq}）：
\[
\bigl(\eu^{\iu\theta}\bigr)^{-m}
= \bigl(\eu^{\iu m\theta}\bigr)^{-1}
= \eu^{\iu(-m)\theta} . \qedhere
\]
\end{proof}

\begin{example}[用棣莫弗公式展开 $\cos 3\theta$]\label{ex:b1:complex:cos3}
记 $c = \cos\theta$，$s = \sin\theta$。由棣莫弗公式与二项式
定理可得
\[
\cos 3\theta + \iu \sin 3\theta = (c + \iu s)^3
= c^3 - 3cs^2 + \iu\,(3c^2 s - s^3),
\]
比较实部，再代入 $s^2 = 1 - c^2$：
\[
\cos 3\theta = c^3 - 3c(1 - c^2) = 4\cos^3\theta - 3\cos\theta .
\]
虚部顺带给出 $\sin 3\theta = 3\sin\theta -
4\sin^3\theta$。一个复恒等式总带有
\emph{两个}实恒等式。反过来读，上面的恒等式是
经典三等分方程的关键：
构造 $\cos(\theta/3)$（从 $\cos\theta$ 出发）等价于求解三次方程 $4x^3 -
3x = \cos\theta$，这正是
在\cref{ch:b1:poly}中的代数理论接手之处。
\end{example}

\begin{proposition}[欧拉公式]\label{prop:b1:complex:euler}
\[
\cos\theta = \frac{\eu^{\iu\theta} + \eu^{-\iu\theta}}{2},
\qquad
\sin\theta = \frac{\eu^{\iu\theta} - \eu^{-\iu\theta}}{2\iu}.
\]
\end{proposition}

\begin{proof}
将 $\eu^{\iu\theta} = \cos\theta +
\iu\sin\theta$ 与 $\eu^{-\iu\theta} = \cos\theta - \iu\sin\theta$ 分别相加、相减。
\end{proof}

\begin{method}[用指数函数处理三角]\label{met:b1:complex:trig}
\begin{enumerate}
  \item \emph{线性化} $\cos^p\theta \sin^q\theta$（把幂化为
        $\cos k\theta$、$\sin k\theta$ 的和）：代入欧拉
        公式，用二项式定理展开
        （\cref{thm:b1:counting:binomial}），再合并共轭项。
  \item \emph{展开} $\cos n\theta$，写成 $\cos\theta$ 的多项式：
        记 $\cos n\theta = \Re\bigl((\cos\theta +
        \iu\sin\theta)^n\bigr)$，展开后将
        偶次幂的 $\sin$ 借助 $\sin^2 = 1 - \cos^2$ 化去。
  \item \emph{三角级数求和}，如 $\sum_k \cos k\theta$：
        识别几何和 $\sum_k
        (\eu^{\iu\theta})^k$ 的实部。
  \item \emph{半角因式分解}：对所有 $p, q$，
        \[
        \eu^{\iu p} + \eu^{\iu q}
        = 2 \cos\tfrac{p - q}{2}\; \eu^{\iu \frac{p+q}{2}},
        \qquad
        \eu^{\iu p} - \eu^{\iu q}
        = 2\iu \sin\tfrac{p - q}{2}\; \eu^{\iu \frac{p+q}{2}} .
        \]
\end{enumerate}
\end{method}

\begin{example}[线性化]\label{ex:b1:complex:linearize}
\[
\cos^3\theta
= \Bigl(\frac{\eu^{\iu\theta} + \eu^{-\iu\theta}}{2}\Bigr)^{\!3}
= \frac{\eu^{3\iu\theta} + 3\eu^{\iu\theta} + 3\eu^{-\iu\theta} +
\eu^{-3\iu\theta}}{8}
= \frac{\cos 3\theta + 3\cos\theta}{4}.
\]
这种形式可以直接积分——这正是线性化
在\cref{ch:b1:integration}中如此重要的缘故。混合乘积的处理方法与此相同，
只是同时用上两个欧拉公式：
\[
\sin^2\theta\cos^2\theta
= \Bigl(\frac{\sin2\theta}2\Bigr)^{\!2}
= \frac{1}{4}\cdot
\Bigl(\frac{\eu^{2\iu\theta} - \eu^{-2\iu\theta}}{2\iu}
\Bigr)^{\!2}
= \frac{2 - \eu^{4\iu\theta} - \eu^{-4\iu\theta}}{16}
= \frac{1 - \cos4\theta}{8} ,
\]
其中第一步用到的二倍角捷径省去了一次
二项展开——在用机械方法计算之前，
总是值得先留意是否存在这类捷径。
\end{example}

\begin{example}[二项三角和]\label{ex:b1:complex:binomsum}
对 $n \in \N$ 与 $\theta \in \R$，计算 $S =
\sum_{k=0}^{n}\binom nk \cos k\theta$。注意它是某个
二项展开式的实部：
\[
S = \Re\sum_{k=0}^n \binom nk \bigl(\eu^{\iu\theta}\bigr)^k
= \Re\bigl(1 + \eu^{\iu\theta}\bigr)^n ,
\]
再对半角作因式分解（见\cref{met:b1:complex:trig}~(4)）：$1 +
\eu^{\iu\theta} = 2\cos\frac\theta2\,\eu^{\iu\theta/2}$，于是
\[
S = \Re\Bigl(2^n\cos^n\frac\theta2\;\eu^{\iu n\theta/2}\Bigr)
= 2^n \cos^n\frac\theta2\,\cos\frac{n\theta}2 .
\]
取虚部即得 $\sum_k\binom nk\sin k\theta =
2^n\cos^n\frac\theta2\sin\frac{n\theta}2$，无需额外计算。验算如下：
取 $\theta = 0$ 可得 $\sum\binom nk = 2^n$，又取 $\theta = \pi$，
则 $S = 0$（$n \geq 1$ 时，每个因子 $\cos\frac\pi2$
为零），即
\cref{ex:b1:counting:binomial}中的交错行和。该方法——“将余弦和看作
复数幂的实部，再对半角作因式分解”——
正是\cref{exo:b1:complex:6}中的方法，只是
以二项式定理代替了等比级数。
\end{example}

\section{复数的方根}

\begin{theorem}[$n$ 次方根]\label{thm:b1:complex:roots}
设 $a = r\,\eu^{\iu\alpha} \neq 0$，$n \in \N^*$。方程
$z^n = a$ 恰有 $n$ 个解：
\[
z_k = r^{1/n}\, \eu^{\iu\left(\frac{\alpha}{n} +
\frac{2k\pi}{n}\right)},
\qquad k = 0, 1, \dots, n - 1 .
\]
\end{theorem}

\begin{proof}
设 $z = \rho\,\eu^{\iu\theta}$（$\rho > 0$）。于是 $z^n =
\rho^n \eu^{\iu n\theta} = r \eu^{\iu\alpha}$ 当且仅当 $\rho^n
= r$（比较模）且 $n\theta \equiv \alpha \pmod{2\pi}$（比较辐角），
即 $\rho = r^{1/n}$ 且 $\theta = \frac{\alpha}{n} +
\frac{2k\pi}{n}$，$k \in \Z$。还需考察何时两个
整数 $k, k'$ 给出同一个数，恰好发生在
两个角相差 $2\pi$ 的整数倍时，
\[
\frac{2k\pi}n - \frac{2k'\pi}n \in 2\pi\Z
\iff \frac{k - k'}n \in \Z
\iff n \mid k - k' .
\]
由带余除法，每个 $k \in \Z$ 模 $n$ 同余于
$\{0, 1, \dots, n-1\}$ 中恰有一个元素，因此这一范围恰好列出
每个解恰出现一次，解的个数恰为 $n$。（同样的
论证在 $\mathbb U_n$ 内重做一遍，可知这些根构成一个正
$n$ 边形：相邻的 $k$ 值所对应的根依次相差固定角
$\frac{2\pi}n$ 的旋转。）
\end{proof}

\begin{example}[$-27$ 的立方根]\label{ex:b1:complex:cuberoots}
解方程 $z^3 = -27$。右端的指数形式：$-27 =
27\,\eu^{\iu\pi}$，故三个根为
\[
z_k = 3\,\eu^{\iu(\frac\pi3 + \frac{2k\pi}3)}, \quad k = 0, 1, 2 :
\qquad
z_0 = 3\eu^{\iu\pi/3} = \frac32 + \frac{3\sqrt3}2\,\iu,
\quad
z_1 = -3,
\quad
z_2 = \conj{z_0} .
\]
下面作两项检验。其一，实根 $-3$ 显然是一个根，而
另外两个根分别是它旋转 $\pm\frac{2\pi}3$ 所得，等价地
即 $-3j$ 与 $-3j^2$。其二，代数上可以验证：$z^3 + 27 = (z +
3)(z^2 - 3z + 9)$，二次式的判别式为 $9 - 36 =
-27 < 0$，其根为 $\frac{3 \pm 3\iu\sqrt3}2 = z_0, \conj{z_0}$。
由此可见：当右端为实数时，非实根总是
成共轭对出现，因此解集的图形
关于实轴对称——这预示了
\cref{ch:b1:poly}中的实因式分解定理。
\end{example}

\begin{example}[右端为非实数的情形]\label{ex:b1:complex:fourthroots}
解方程 $z^4 = -8 + 8\iu\sqrt3$。右端的指数形式：
模为 $\sqrt{64 + 192} = 16$，辐角 $\theta$ 满足
$\cos\theta = -\frac12$，$\sin\theta = \frac{\sqrt3}2$，即
$\theta = \frac{2\pi}3$。四个根为
\[
z_k = 2\,\eu^{\iu(\frac\pi6 + \frac{k\pi}2)}, \quad k = 0, 1, 2,
3 :
\qquad
z_0 = \sqrt3 + \iu,\quad z_1 = \iu z_0 = -1 + \iu\sqrt3,
\]
\[
z_2 = -z_0 = -\sqrt3 - \iu,\qquad z_3 = -\iu z_0 = 1 - \iu\sqrt3 .
\]
（验算：$z_0^2 = 2 + 2\iu\sqrt3$，于是 $z_0^4 = (2 +
2\iu\sqrt3)^2 = 4 - 12 + 8\iu\sqrt3 = -8 + 8\iu\sqrt3$。）
一旦求得\emph{一个}根，其余三个根便随之得到：它们是该根
依次旋转 $\frac\pi2$ 所得，即它与四次单位根的乘积——
这正是
\cref{thm:b1:complex:roots}背后的一般结构，值得在从头逐个重算之前加以利用。
这里没有共轭对称性：
右端并非实数。
\end{example}

\begin{definition}[单位根]\label{def:b1:complex:unity}
\emph{$n$ 次单位根}\index{单位根}是
方程 $z^n = 1$ 的解：
\[
\mathbb{U}_n = \bigl\{\, \omega^k : k = 0, \dots, n-1 \,\bigr\},
\qquad \omega = \eu^{2\iu\pi/n}.
\]
它们在乘法下构成群（见\cref{ch:b1:structures}），
且位于单位圆内接正 $n$ 边形的顶点上。
\end{definition}

\begin{omfigure}
\begin{tikzpicture}[scale=1.5]
  \draw[->, thin] (-1.35,0) -- (1.45,0) node[below] {$\Re$};
  \draw[->, thin] (0,-1.35) -- (0,1.45) node[left] {$\Im$};
  \draw[thin, gray] (0,0) circle (1);
  \foreach \k in {0,...,4} {
    \fill[omDef] ({cos(72*\k)},{sin(72*\k)}) circle (0.035);
  }
  \draw[omDef, thick] (1,0) -- ({cos(72)},{sin(72)}) -- ({cos(144)},{sin(144)})
    -- ({cos(216)},{sin(216)}) -- ({cos(288)},{sin(288)}) -- cycle;
  \node[right, font=\small, omDef] at (1.02,0.12) {$1$};
  \node[above right, font=\small, omDef] at ({cos(72)},{sin(72)})
    {$\omega$};
  \node[above left, font=\small, omDef] at ({cos(144)},{sin(144)})
    {$\omega^2$};
  \node[below left, font=\small, omDef] at ({cos(216)},{sin(216)})
    {$\omega^3$};
  \node[below right, font=\small, omDef] at ({cos(288)},{sin(288)})
    {$\omega^4$};
\end{tikzpicture}

{\small 五次单位根，$\omega = \eu^{2\iu\pi/5}$：单位圆上的
正五边形。}
\end{omfigure}

\begin{proposition}[单位根之和]\label{prop:b1:complex:sumroots}
当 $n \geq 2$ 时，$n$ 次单位根之和为零：
$\sum_{k=0}^{n-1} \omega^k = 0$。
\end{proposition}

\begin{proof}
公比 $\omega \neq 1$ 的等比求和：
$\sum_{k=0}^{n-1} \omega^k = \frac{\omega^n - 1}{\omega - 1} = 0$
由于 $\omega^n = 1$。
\end{proof}

\begin{example}[读取实部与虚部]\label{ex:b1:complex:sumrootsparts}
将 $\sum_{k=0}^{n-1}\omega^k = 0$ 分为实部与
虚部，可直接得到两个三角恒等式：
\[
\sum_{k=0}^{n-1}\cos\frac{2k\pi}n = 0,
\qquad
\sum_{k=0}^{n-1}\sin\frac{2k\pi}n = 0
\qquad (n \geq 2).
\]
从几何上看，单位圆内接正 $n$ 边形的重心
即圆心 --- 各顶点恰好达到平衡。
当 $n = 5$ 时，第一个恒等式给出 $1 + 2\cos\frac{2\pi}5 +
2\cos\frac{4\pi}5 = 0$（将 $k$ 与 $n - k$ 配对），这正是
\cref{exo:b1:complex:8}中计算
$\cos\frac{2\pi}5$ 的出发点。
\end{example}

\begin{example}[代数形式的平方根]\label{ex:b1:complex:sqrt}
不借助三角函数，求解 $z^2 = 3 + 4\iu$，设 $z = x + \iu y$：
\[
x^2 - y^2 = 3, \qquad 2xy = 4, \qquad x^2 + y^2 = \abs{3 + 4\iu} = 5 .
\]
将第一式与第三式相加得 $x^2 = 4$，于是 $x = \pm 2$，再由 $y = 2/x =
\pm 1$，符号搭配须\emph{相同}（$xy = 2 > 0$）：$z = \pm(2 +
\iu)$。结合通常的求根公式，即可解出
任意带复系数的二次方程
（见\cref{exo:b1:complex:7}）。
\end{example}

\section{复数与平面几何}

\begin{proposition}[几何对照]\label{prop:b1:complex:geometry}
将平面上的点 $M(x, y)$ 与其\emph{复坐标} $z = x
+ \iu y$\index{复坐标}。设 $A, B, C$ 为两两不同的点，其复坐标依次为 $a, b,
c$：
\begin{enumerate}
  \item $\abs{b - a}$ 是 $AB$ 的距离；
  \item $\arg\dfrac{c - a}{b - a}$ 是向量
        $\vect{AB}$ 与 $\vect{AC}$ 之间的夹角（模 $2\pi$）；
  \item $A, B, C$ 共线当且仅当 $\dfrac{c - a}{b - a} \in
        \R$；直线 $AB$ 与 $AC$ 垂直当且仅当
        $\dfrac{c - a}{b - a} \in \iu\R$。
\end{enumerate}
\end{proposition}

\begin{proof}
（1）是将模的定义用于 $b - a$，它是
$\vect{AB}$ 的复坐标。（2）：记 $b - a = r\eu^{\iu\theta}$，$c - a =
s\eu^{\iu\varphi}$；则 $\frac{c-a}{b-a} = \frac sr
\eu^{\iu(\varphi - \theta)}$ 的辐角为 $\varphi - \theta$，即
由 $\vect{AB}$ 到 $\vect{AC}$ 的夹角。（3）：共线即夹角为 $0$
或 $\pi$，即辐角属于 $\pi\Z$，亦即商为实数；
垂直即夹角为 $\pm\frac\pi2$，亦即商
为纯虚数。（由 $C \neq A$，商非零。）
\end{proof}

\begin{remark}[插曲：$\C$ 是赋予了乘法的平面]\label{rem:b1:complex:interlude}
值得停下来想一想，本章究竟依赖什么才得以成立：
平面 $\R^2$ 在各个方向上都有加法，
却没有天然的乘法 --- 而 $\C$ \emph{正是}
装备了这样一种乘法的平面，在其中，乘以一个固定的数
即旋转并伸缩。这一结构在本卷中还将
被反复运用三次。在\cref{ch:b1:matrices}中，乘以
$a + \iu b$ 再现为 $2 \times 2$ 矩阵，其两行为 $(a,
-b)$ 与 $(b, a)$：复数运算是第一族
完全具体的矩阵乘积。在\cref{ch:b1:euclid}中，
公式 $\Re(\conj z\,w)$ 恰为内积，
而 $\abs z$ 即欧几里得范数：本章证明的三角不等式
正是那里柯西–施瓦茨不等式的原型。在
\cref{ch:b1:curves}中，动点最好写作 $t \mapsto
z(t)$，使速度与加速度成为复值
导数 --- 例如圆周运动，就是简单地 $z(t) =
R\,\eu^{\iu\omega t}$。一种好的乘法，带来四章的
收益。
\end{remark}

\begin{proposition}[映射 $z \mapsto az + b$]\label{prop:b1:complex:similitude}
设 $a \in \C^*$，$b \in \C$。平面的变换 $f(z) = az + b$
有如下性质：
\begin{itemize}
  \item 当 $a = 1$ 时为平移；
  \item 否则有唯一不动点 $\zeta = \frac{b}{1-a}$，
        且 $f(z) - \zeta = a\,(z - \zeta)$：$f$ 是
        中心为 $\zeta$、角为 $\arg a$ 的旋转，与
        中心为 $\zeta$、比例为 $\abs a$ 的位似变换的复合。
\end{itemize}
特别地，$z \mapsto \eu^{\iu\theta} z$ 是角为
$\theta$ 的绕原点旋转，$\conj z$ 是关于实
轴的反射。
\end{proposition}

\begin{proof}
若 $a = 1$，$f(z) = z + b$ 是沿复数 $b$ 所对应的向量平移。若
$a \neq 1$，则不动点方程 $z = az + b$ 有唯一
解 $\zeta = \frac{b}{1-a}$，于是 $f(z) - \zeta = az + b -
(a\zeta + b) = a(z - \zeta)$。记 $a = \abs a\, \eu^{\iu\arg a}$，
乘以 $a$ 使到 $\zeta$ 的距离变为原来的 $\abs a$ 倍，并
使各角增加 $\arg a$，各角顶点均为 $\zeta$，这正是前面所述的复合。
\end{proof}

\begin{example}[判别映射 $z \mapsto az+b$]\label{ex:b1:complex:classify}
取 $f(z) = \iu z + 1$。此时 $a = \iu \neq 1$，不动点为
\[
\zeta = \frac{b}{1 - a} = \frac1{1 - \iu}
= \frac{1 + \iu}{2},
\]
由于 $\abs a = 1$，完全没有伸缩：$f$ 是
中心为 $\frac{1+\iu}2$、角为 $\arg \iu = \frac\pi2$ 的纯旋转。
验证：$f(\zeta) = \iu\,\frac{1+\iu}2 + 1 = \frac{\iu - 1}2 + 1 =
\frac{1 + \iu}2 = \zeta$，且 $f(0) = 1$，$f(1) = 1 + \iu$，$f(1 +
\iu) = \iu\,(1 + \iu) + 1 = \iu$：四点 $0, 1, 1+\iu, \iu$
是单位正方形的顶点，绕中心 $\zeta$ 循环，
每次转过四分之一圈 --- 恰是
绕正方形中心作 $\frac\pi2$ 旋转所应做到的。
\end{example}

\begin{omfigure}
\begin{tikzpicture}[scale=1.6]
  \draw[->, thin] (-0.5,0) -- (1.6,0) node[below] {$\Re$};
  \draw[->, thin] (0,-0.35) -- (0,1.5) node[left] {$\Im$};
  \draw[gray, thin] (0,0) rectangle (1,1);
  \fill[omDef] (0.5,0.5) circle (0.03)
    node[below right, font=\small] {$\zeta$};
  \fill (0,0) circle (0.025) node[below left, font=\small] {$0$};
  \fill (1,0) circle (0.025) node[below right, font=\small] {$1$};
  \fill (1,1) circle (0.025)
    node[above right, font=\small] {$1{+}\iu$};
  \fill (0,1) circle (0.025) node[above left, font=\small] {$\iu$};
  \draw[omThm, ->, thick] (0.08,-0.12)
    .. controls (0.6,-0.28) and (1.12,0.4) .. (1.12,0.9);
  \draw[omThm, ->, thick] (1.12,1.08)
    .. controls (0.9,1.3) and (0.35,1.3) .. (0.05,1.12);
  \node[omThm, font=\small] at (0.95,-0.28) {$0 \mapsto 1$};
  \node[omThm, font=\small] at (1.62,0.75) {$1 \mapsto 1{+}\iu$};
\end{tikzpicture}

{\small 映射 $f(z) = \iu z + 1$ 出自
\cref{ex:b1:complex:classify}：绕不动点
$\zeta = \frac{1+\iu}2$ 的四分之一圈旋转。单位正方形的
顶点按 $0 \to 1 \to 1{+}\iu \to \iu \to 0$ 循环；没有任何点沿直线
运动，但整个正方形被刚性地转动。}
\end{omfigure}

\begin{remark}[关于模与辐角的常见误区]\label{rem:b1:complex:pitfalls}
\begin{enumerate}
  \item \emph{$\C$ 中不存在不等式。}对非实数写 $z \leq w$
        的做法没有意义；只有模、实部
        与虚部可以比较。
  \item \emph{$\sqrt{\phantom z}$ 专用于 $\R_+$。}每个
        非零复数有\emph{两个}平方根，二者
        无优先之分：写“设 $\delta$ 为 $\Delta$ 的一个
        平方根”（计算方法同
        \cref{ex:b1:complex:sqrt}），切勿写 $\sqrt\Delta$ ---
        规则 $\sqrt{ab} = \sqrt a\sqrt b$ 在
        $a = b = -1$ 时即已不成立。
  \item \emph{辐角按模 $2\pi$ 理解。}由
        $\eu^{\iu\alpha} = \eu^{\iu\beta}$，可得 $\alpha
        \equiv \beta \pmod{2\pi}$，而非 $\alpha = \beta$；
        忘记 $2k\pi$ 会使方程 $z^n = a$ 的解集
        丢失 $n - 1$ 个（共 $n$ 个）元素。
  \item \emph{$\abs{z + w}$ 不等于 $\abs z + \abs w$。}
        三角不等式中的等号成立，是例外的共线
        情形（见\cref{prop:b1:complex:rules}~(4)）；一般地，
        和的模只能估计，不能直接算出。
\end{enumerate}
\end{remark}

\begin{example}\label{ex:b1:complex:equilateral}
$A, B, C$（对应复数 $a, b, c$，两两不同）构成一个等边
三角形，顶点按正向（逆时针）排列，当且仅当
$\frac{c - a}{b - a} = \eu^{\iu\pi/3}$：中心为
$A$、角为 $\frac\pi3$ 的旋转把 $B$ 映为 $C$。两种
定向合起来，可由对称方程 $a^2 + b^2 + c^2 = ab
+ bc + ca$ 统一刻画（见\cref{exo:b1:complex:10}）。
\end{example}

\begin{remark}[本章的应用]\label{rem:b1:complex:whereused}
指数形式是本卷中复用最频繁的
计算工具。在
\cref{ch:b1:structures}中，单位根成为循环群的标准例子，并
在\cref{ch:b1:poly}中推动 $X^n - 1$ 的因式分解；
将其中的共轭根成组，即得在
\cref{ch:b1:fractions}中部分分式所用的实因式分解。线性化（见\cref{met:b1:complex:trig}）
是积分三角幂的标准准备，在
\cref{ch:b1:integration}中；而\cref{ch:b1:diffeq}的特征方程
有复根，其实部与虚部
给出振荡解 $\eu^{\lambda t}\cos\omega
t$。几何的对应以矩阵的形式再次出现：旋转与
相似变换化为
\cref{ch:b1:matrices,ch:b1:euclid}中的正交矩阵，并且在
\cref{ch:b1:curves}中的参数曲线常以映射 $t \mapsto z(t)
\in \C$ 的形式表示最为方便。
\end{remark}

\section{练习}

\begin{exercise}[$\star$]\label{exo:b1:complex:1}
将下列各数化为指数形式：$1 + \iu$；$\;\sqrt 3 - \iu$；$\;-5$；
$\;\dfrac{1 + \iu}{\sqrt 3 - \iu}$。由此推出
$\cos\frac{5\pi}{12}$ 与 $\sin\frac{5\pi}{12}$ 的值。
\end{exercise}

\begin{exercise}[$\star$]\label{exo:b1:complex:2}
计算 $(1 + \iu)^{20}$。对于哪些 $n \in \N$，$(1 + \iu)^n$ 是
实数？
\end{exercise}

\begin{exercise}[$\star$]\label{exo:b1:complex:3}
在几何上描述满足下列条件的 $z \in \C$ 构成的集合：
$\abs{z - 2} = \abs{z + \iu}$；$\;\abs{z - 1} = 2$；
$\;\dfrac{z - 1}{z + 1} \in \iu\R$（$z \neq -1$）。
\end{exercise}

\begin{exercise}[$\star$]\label{exo:b1:complex:4}
将 $\sin^4\theta$ 线性化，并把 $\cos 4\theta$ 展开为关于
$\cos\theta$ 的多项式。
\end{exercise}

\begin{exercise}[$\star$]\label{exo:b1:complex:5}
解方程 $z^3 = 8\iu$，并把解画在图上。解方程
$z^4 = -4$，并将 $X^4 + 4$ 分解为两个实二次多项式之积。
\end{exercise}

\begin{exercise}[$\star\star$]\label{exo:b1:complex:6}
设 $\theta \in \R$ 满足 $\eu^{\iu\theta} \neq 1$，且 $n \in \N$，
计算
\[
C_n = \sum_{k=0}^{n} \cos k\theta
\qquad\text{和}\qquad
S_n = \sum_{k=0}^{n} \sin k\theta
\]
通过对几何级数 $\sum_k \eu^{\iu k\theta}$ 求和，并利用
半角因式分解。
\end{exercise}

\begin{exercise}[$\star\star$]\label{exo:b1:complex:7}
在 $\C$ 中解方程：$z^2 - (3 + 4\iu) z + (-1 + 5\iu) = 0$。\emph{（计算
判别式，并按
\cref{ex:b1:complex:sqrt}中的方法求平方根。）}
\end{exercise}

\begin{exercise}[$\star\star$]\label{exo:b1:complex:8}
设 $\omega = \eu^{2\iu\pi/5}$。
\begin{enumerate}
  \item 证明 $1 + \omega + \omega^2 + \omega^3 + \omega^4 = 0$。
  \item 令 $u = \omega + \omega^4$ 与 $v = \omega^2 + \omega^3$。
        计算 $u + v$ 与 $uv$，由此推出 $u$ 与 $v$ 是
        $X^2 + X - 1$ 的根。
  \item 由此得 $\cos\frac{2\pi}{5} = \frac{\sqrt 5 - 1}{4}$。
\end{enumerate}
\end{exercise}

\begin{exercise}[$\star\star$]\label{exo:b1:complex:9}
证明：对任意 $z, w \in \C$，有（平行四边形恒等式）：
\[
\abs{z + w}^2 + \abs{z - w}^2 = 2\abs z^2 + 2\abs w^2 ,
\]
并在以 $0,
z, w, z + w$ 为顶点的平行四边形中给出几何解释。
\end{exercise}

\begin{exercise}[$\star\star\star$]\label{exo:b1:complex:10}
证明：复坐标分别为 $a, b, c$ 的三个两两不同的点构成
正三角形（不论定向）当且仅当
\[
a^2 + b^2 + c^2 = ab + bc + ca .
\]
\emph{提示：直接情形为 $\frac{c-a}{b-a} = -j^2$，而
间接情形为 $\frac{c-a}{b-a} = -j$，其中 $j = \eu^{2\iu\pi/3}$
满足 $j^2 + j + 1 = 0$；或者分解 $a + jb + j^2c$ 与 $a + j^2 b
+ jc$。}
\end{exercise}

\begin{exercise}[$\star\star\star$]\label{exo:b1:complex:11}
对 $n \in \N^*$，求 $P = \prod_{k=1}^{n-1}
\bigl(1 - \omega^k\bigr)$ 的值，其中 $\omega = \eu^{2\iu\pi/n}$。
\emph{提示：$X^n - 1 = \prod_{k=0}^{n-1} (X - \omega^k)$；两边除以
$X - 1$ 后代入 $X = 1$。} 由此推出
$\prod_{k=1}^{n-1} \sin\frac{k\pi}{n} = \dfrac{n}{2^{n-1}}$。
\end{exercise}

\begin{exercise}[$\star\star$]\label{exo:b1:complex:12}
设 $n \geq 2$。解方程 $(z + 1)^n = (z - 1)^n$（在 $\C$ 中）：
证明它恰有 $n - 1$ 个解，且均为纯虚数，
即
\[
z_k = -\iu\,\frac{\cos\frac{k\pi}{n}}{\sin\frac{k\pi}{n}},
\qquad k = 1, \dots, n - 1 .
\]
\emph{提示：$z = 1$ 不是解，因此可作除法，并利用
单位根；再应用
\cref{met:b1:complex:trig}中的半角因式分解。}
\end{exercise}

\section{问题：拿破仑定理}

\begin{problem}\label{pb:b1:complex:1}
在一个
\emph{任意}三角形的每条边外侧各作一个正三角形，这三个正三角形的中心
总构成一个正三角形。这就是拿破仑定理
\index{拿破仑定理} --- 一个看不出任何成立理由的命题，
其成立性由 $j = \eu^{2\iu\pi/3}$ 的代数在三行计算中证明。
本题将建立完整的工具：
（旋转、直相似，以及 $j$ 判别准则，
见\cref{exo:b1:complex:10}），证明拿破仑定理的内、
外两种情形，指出退化情形，并以同一学派的第二颗明珠作结：
即关于正三角形的范斯霍滕定理，
以及托勒密不等式。全文中，平面上的点
均等同于其复坐标。

\medskip
\noindent\textbf{第一部分 --- 数 $j$ 与旋转。}
\begin{enumerate}
  \item 求 $j$ 的代数形式，再求 $j^2$、$j^3$、$1 + j
        + j^2$、$\conj j$ 以及 $j^{-1}$，并在单位圆的
        示意图上标出 $1$、$j$、$j^2$。
  \item 证明中心为 $a$、旋转角为 $\theta$ 的旋转为
        $r(z) = a + \eu^{\iu\theta}(z - a)$，并证明
        两个旋转（旋转角分别为 $\theta$ 与 $\theta'$）的复合，
        是旋转角为 $\theta + \theta'$ 的旋转
        （当 $\theta + \theta' \notin 2\pi\Z$ 时），否则为平移。
        （利用\cref{prop:b1:complex:similitude}。）
  \item 设 $b \neq c$。证明恰有两个点 $p$
        使 $(b, c, p)$ 为正三角形，它们为 $p = b +
        \eu^{\pm\iu\pi/3}(c - b)$。对于\emph{定向}三角形
        $(a, b, c)$，在例子 $a = 0$、$b = 1$、$c =
        \iu$ 中验证：选取 $\eu^{-\iu\pi/3}$ 的那个顶点，
        恰位于直线 $BC$ 上与 $a$ 相对的一侧，即
        \emph{外侧}顶点。
  \item 由\cref{exo:b1:complex:10}的解答，$(a, b, c)$
        为定向正三角形当且仅当 $a + jb + j^2c = 0$。证明下面两个
        旋转恒等式
        \[
        b + jc + j^2a = j^2\,(a + jb + j^2c),
        \qquad
        c + ja + j^2b = j\,(a + jb + j^2c),
        \]
        并由此推出该判别准则在 $(a, b, c)$ 的
        循环置换下不变。
  \item 证明若 $a + jb + j^2c = 0$ 且三点中有两点
        重合，则三点全部重合。（因此该判别准则
        恰好刻画了：定向正三角形，
        或单个点。）
\end{enumerate}

\medskip
\noindent\textbf{第二部分 --- 直相似。}
\begin{enumerate}[resume]
  \item 证明形如 $f(z) = az + b$（$a \in \C^*$）的映射
        （\emph{直相似}）对复合与求逆封闭：
        其中两个的复合，以及其中任一个的逆，
        仍具有这种形式。（用
        \cref{ch:b1:structures}中的语言说，它们构成一个群。）
  \item 给定 $z_1 \neq z_2$ 与 $w_1 \neq w_2$，证明存在
        \emph{唯一}的直相似 $f$，使 $f(z_1) = w_1$ 且
        $f(z_2) = w_2$，并显式给出 $a$ 与 $b$。
  \item 证明直相似 $f(z) = az + b$ 将所有距离乘以 $\abs a$，并保持
        任意三角形的\emph{形状}
        $\frac{c - a'}{b' - a'}$ 不变，此处 $(a', b', c')$ 为任一三角形
        \emph{（分子与分母都乘以
        $a$）}。由此推出两个三角形直相似
        当且仅当它们的形状相等。
  \item 完全确定满足 $f(0) = 1$
        与 $f(1) = \iu$ 的直相似：给出 $a$、$b$、不动点、比值
        与旋转角。
  \item 设 $\alpha + \beta = 1$。证明“加权点”
        $g(a', b') = \alpha a' + \beta b'$ 与每个直相似可交换：
        $f(\alpha a' + \beta b') = \alpha f(a')
        + \beta f(b')$。由此推出重心、中点以及下文的
        拿破仑中心都被相似变换保持
        --- 这正是每一个“不失一般性，设外接圆
        为单位圆”论证背后的代数依据。
        论证背后的代数依据。
\end{enumerate}

\medskip
\noindent\textbf{第三部分 --- 拿破仑定理。}设 $(a, b, c)$
        为定向三角形。在每条边外侧各作一个正三角形
        （见第 3 题），并设 $n_a$、$n_b$、$n_c$ 分别为
        在 $[b, c]$、$[c, a]$、$[a, b]$ 上所作三角形的
        中心（重心）。
\begin{enumerate}[resume]
  \item 证明
        \[
        n_a = \alpha b + \beta c, \qquad
        n_b = \alpha c + \beta a, \qquad
        n_c = \alpha a + \beta b,
        \qquad\text{其中}\quad
        \alpha = \frac{3 + \iu\sqrt3}{6},\ \beta = \conj\alpha .
        \]
  \item 验证 $\alpha + \beta = 1$，并由此推出三角形
        $(n_a, n_b, n_c)$ 与
        $(a, b, c)$ 具有\emph{相同的重心}。
  \item 利用第 4 题中的恒等式，证明
        \[
        n_a + j\,n_b + j^2 n_c
        = (\alpha j^2 + \beta j)\,(a + jb + j^2c) .
        \]
  \item 计算 $\alpha j + \beta$，并作出结论：$\alpha j^2 +
        \beta j = j(\alpha j + \beta) = 0$，故 $n_a + j n_b + j^2
        n_c = 0$，对\emph{每一个}三角形都成立：外侧
        拿破仑三角形为定向正三角形（或为一个点）。
        拿破仑定理得证。
  \item 改为向\emph{内侧}作正三角形
        （取第 3 题中的 $\eu^{+\iu\pi/3}$），并设 $m_a, m_b,
        m_c$ 为它们的中心。证明 $m_a = \beta b + \alpha c$（其余
        类推），再证明 $m_a + j^2 m_b + j\,m_c = 0$：内侧
        拿破仑三角形同样为正三角形，
        但定向相反。
  \item 找出退化情形：证明 $n_a = n_b = n_c$
        当且仅当 $a + j^2 b + jc = 0$，即 $(a, b, c)$
        为\emph{反向}正三角形。\emph{（
        $n_a, n_b, n_c$ 中两点重合当且仅当三点全部重合，由第 5 题可得；
        然后利用 $n_a + j^2 n_b + j n_c = j(\alpha +
        \beta j)(a + j^2b + jc)$，并验证 $\alpha + \beta j \neq
        0$。）}
  \item 对三角形 $a = 0$、
        $b = 1$、$c = \iu$ 完成全部计算：精确给出 $n_a, n_b, n_c$，并
        通过直接计算三条边长的平方来验证该三角形
        为正三角形，其边长平方为
        $\frac{2 + \sqrt3}{3}$。
\end{enumerate}

\medskip
\noindent\textbf{第 IV 部分 --- 托勒密与范斯霍滕。}
\begin{enumerate}[resume]
  \item 证明对所有复数 $a, b, c, d$ 成立的恒等式：
        \[
        (a - b)(c - d) + (a - d)(b - c) = (a - c)(b - d) .
        \]
  \item 由此推出\emph{托勒密不等式}\index{托勒密不等式}：
        对任意四点 $A, B, C, D$，
        \[
        AC \cdot BD \;\leq\; AB \cdot CD + AD \cdot BC ,
        \]
        等号成立，当且仅当 $(a-b)(c-d)$ 与
        $(a-d)(b-c)$ 位于从 $0$ 出发的同一条射线上。
  \item 设 $\theta, \varphi \in \R$，证明
        $\abs{\eu^{\iu\theta} - \eu^{\iu\varphi}} =
        2\,\abs{\sin\frac{\theta - \varphi}2}$。
  \item （范斯霍滕定理\index{范斯霍滕定理}）设
        $(a, b, c) = (1, j, j^2)$ ---由第 II 部分可知，此设定
        对正向等边三角形不失一般性 ---并取
        $p = \eu^{\iu\theta}$，其中 $\theta \in
        \intoo{2\pi/3}{4\pi/3}$，它是外接圆上
        位于弧 $BC$ 上且不含 $A$ 的点，证明
        \[
        PA = PB + PC .
        \]
        \emph{（用第 20 题表示三个距离，并利用
        和差化积公式。）}
  \item 验证当 $p = -1$ 时，这恰好是
        对循环顺序 $A, B, P, C$ 的托勒密不等式取等号的情形：
        计算 $(a - b)(p - c)$ 与 $(a - c)(b - p)$，并验证
        它们的比为正实数。
\end{enumerate}

\medskip
\noindent\textbf{第 V 部分 --- 综合。}
\begin{enumerate}[resume]
  \item 计算等边三角形
        $(1, j, j^2)$ 本身的外拿破仑三角形，
        并从几何上描述结果。
  \item 本题究竟在何处用到了：(i) 把乘法看作旋转；
        (ii) 相似变换的群结构，以及
        形状不变量；(iii) 在
        \cref{met:b1:complex:trig} 中的半角分解？每项各用一句话回答。
  \item 用一小段文字阐明本题的寓意：
        平面几何与 $\C$ 的代数之间的“词典”带来什么，
        付出什么代价，以及下面两个证明步骤——
        恒等式 $\alpha j + \beta = 0$ 与经典图形——
        哪一个在你看来更能\emph{解释}拿破仑定理？
        请指出一处，说明这一“词典”
        将以矩阵形式在本卷后文中再次出现
        的地方。
\end{enumerate}
\end{problem}
